\section{Contractions based on Fourier Transforms}

\red{Walsh-Hadamard transform and further transforms can be efficiently implemented by quantum circuits.
We here investigate how these transforms can be utilized to compute contractions of tensor networks.}


\subsection{Walsh-Hadamard transform}

The application of Hadamard gates on the distributed variables afterwards results in a Walsh-Hadamard transform of the prepared state.

The transform is performed by applying Hadamard gates on each variable.
This is equivalent with contraction of the tensor
\begin{align*}
    \hgate^{\catorder}\left[\catvariableof{[\catorder],\insymbol},\catvariableof{[\catorder],\outsymbol}\right] =
    \bigotimes_{\catenumeratorin} \hgateat{\ccatvariableof{\catenumerator}{\insymbol},\ccatvariableof{\catenumerator}{\outsymbol}} \, ,
\end{align*}
which coordinates are
\begin{align*}
    \hgate^{\catorder}\left[\indexedcatvariableof{[\catorder],\insymbol},\indexedcatvariableof{[\catorder],\outsymbol}\right] =
    \sqrt{\frac{1}{2^{\catorder}}} \left(-1\right)^{\sum_{\catenumeratorin}\catindexof{\catenumerator,\insymbol} \cdot \catindexof{\catenumerator,\outsymbol}} \, .
\end{align*}

We have
\begin{align*}
    \contractionof{\onehotmapofat{0}{\ccatvariable{\insymbol}},
        \hgateat{\ccatvariable{\insymbol},\ccatvariable{\outsymbol}}}{\ccatvariable{\outsymbol}}
    &= \sqrt{\frac{1}{2}} \onesat{\ccatvariable{\outsymbol}} \\
    \contractionof{\onehotmapofat{1}{\ccatvariable{\insymbol}},
        \hgateat{\ccatvariable{\insymbol},\ccatvariable{\outsymbol}}}{\ccatvariable{\outsymbol}}
    &= \sqrt{\frac{1}{2}} \coloredmatrixof{1 \\ -1}{\ccatvariable{\outsymbol}}
\end{align*}
and conversely
\begin{align*}
    \contractionof{\sqrt{\frac{1}{2}} \onesat{\ccatvariable{\insymbol}},
        \hgateat{\ccatvariable{\insymbol},\ccatvariable{\outsymbol}}}{\ccatvariable{\outsymbol}}
    &= \onehotmapofat{0}{\ccatvariable{\outsymbol}} \\
    \contractionof{\sqrt{\frac{1}{2}} \coloredmatrixof{1 \\ -1}{\ccatvariable{\insymbol}},
        \hgateat{\ccatvariable{\insymbol},\ccatvariable{\outsymbol}}}{\ccatvariable{\outsymbol}}
    &= \onehotmapofat{1}{\ccatvariable{\outsymbol}} \, .
\end{align*}
If and only if the function is constant, then the transformed state is the ground state.

\begin{remark}[Fourier transform]
    The Walsh-Hadamard transform is the Fourier transform on the group $Z_2^\catorder$, the $\catorder$-ary product of the cyclic group of order $2$.
    The irreducible representations of this group are one-dimensional and their characters are given by the slices of the Hadamard tensor.
\end{remark}



\subsection{Examples}

Let us present two known algorithms as the measurement of states prepared by \computationCircuits{}.

\input{partI/5-quantum-circuits/examples/transform-based/deutsch-jozsa}

Deutsch-Jozsa can be generalized to estimating the moments of a distribution given its q-sample.
Note, that we have to demand a q-sample (i.e. vanishing phase tensor), otherwise the measurement is affected by the phase tensor and the Deutsch-Jozsa test is effectively done wrt another function.
If for example the coordinates of the phase tensor are multiples of $\pi$, and the phases are thus $+1/-1$, then the test is done wrt the function $\exformula\oplus(\frac{1}{\pi}\phasecore)$ instead of $\exformula$.



\input{partI/5-quantum-circuits/examples/transform-based/inversion-test}

%\subsection{Further Overlap-measuring Circuits}
Quantum Circuits such as the SWAP test and the Hadamard test (\cite[Section 3.6.1]{schuld_machine_2021}) can be used to measure overlaps of quantum states, which are the squared absolutes of contractions of two state tensors.
\begin{itemize}
    \item When we have preparation schemes for two tensors, we can control them with a common ancilla qubit and measure their contraction by a Hadamard test (alternatively, using the SWAP test and state preparation in two registers).
    \item Can we extend these schemes to contractions of more general tensor networks?
\end{itemize}

%\subsection{Discussion}
While the preparation of contracted basis encodings is done efficiently, accessing the prepared state by measurements is a bottleneck.
This is a typical \emph{quantum state tomography} problem \cite{roth_semi-device-dependent_2023}.
Sparsity of the contracted state can be exploited to probabilistically guarantee precise estimates \cite{gross_quantum_2010}.
However, when deciding whether $\exformula$ is satisfiable, based on the prepared contraction $\sqrt{\normalizationof{\bencodingofat{\exformula}{\headvariable,\shortcatvariables}}{\headvariable}}$, we have an exponentially small error tolerance.
Moreover, the applied sample mean is the UMVUE (Uniformly Minimum Variance Unbiased Estimator, see \cite{casella_statistical_2001}), and can thus not be improved by other unbiased estimators.


\subsection{Generic Hidden Subgroup Problems}

In general \cite{nielsen_quantum_2010} approached by generic Fourier transforms of the distributed variables of a basis encoding state to a function.